% ==================================================================
%  Preamble: packages, styling, and shared macros.
%  Document metadata (title, version, authors) lives in main.tex and
%  must be defined *before* this file is input.
%  Requires only standard packages (TeX Live base + recommended).
% ==================================================================

% --- Page geometry and fonts --------------------------------------
\usepackage[margin=2.5cm]{geometry}
\usepackage[T1]{fontenc}
\usepackage[utf8]{inputenc}
\usepackage{lmodern}
\usepackage{microtype}

% --- Mathematics ---------------------------------------------------
\usepackage{amsmath}
\usepackage{amssymb}
\usepackage{amsthm}

% Theorem-like environments. They share a single counter, numbered
% within the section (Definition 2.1, Theorem 2.2, ...), so a reference
% is unambiguous whichever kind it points at.
% To add another kind, copy a line and keep the [theorem] counter:
%   \newtheorem{lemma}[theorem]{Lemma}
\theoremstyle{plain}         % bold head, italic body
\newtheorem{theorem}{Theorem}[section]
\newtheorem{proposition}[theorem]{Proposition}
\theoremstyle{definition}    % bold head, upright body
\newtheorem{definition}[theorem]{Definition}
\theoremstyle{plain}         % restore, so later \newtheorem lines inherit it
% amsthm also supplies an unnumbered `proof' environment.

% --- Figures, tables, lists ----------------------------------------
\usepackage{graphicx}
\usepackage{float}
\usepackage{booktabs}
\usepackage{tabularx}
\usepackage{longtable}
\usepackage{array}
\usepackage[font=small,labelfont=bf]{caption}
\usepackage{enumitem}
\graphicspath{{figures/}{./}}

% --- Colour, links, citations --------------------------------------
\usepackage{xcolor}
\usepackage[square,numbers,sort]{natbib}
\definecolor{citecol}{RGB}{0,90,160}
\definecolor{rulecol}{RGB}{0,90,160}
\definecolor{placeholdercol}{gray}{0.45}
\usepackage[colorlinks=true,urlcolor=blue!60!black,linkcolor=black,citecolor=citecol]{hyperref}
\hypersetup{
  pdftitle={\doctitle},
  pdfauthor={\docauthors},
  pdfsubject={\projectname\ --- \projectunit},
  pdfkeywords={cluster analysis, unsupervised learning, \projectname},
  bookmarksnumbered=true,
}

% In-text citations rendered as bold coloured [n] (brackets included).
\NewCommandCopy{\origcitep}{\citep}
\NewCommandCopy{\origcitet}{\citet}
\RenewDocumentCommand{\citep}{oom}{{\bfseries\color{citecol}%
  \IfNoValueTF{#1}{\origcitep{#3}}{\IfNoValueTF{#2}{\origcitep[#1]{#3}}{\origcitep[#1][#2]{#3}}}}}
\RenewDocumentCommand{\citet}{oom}{{\bfseries\color{citecol}%
  \IfNoValueTF{#1}{\origcitet{#3}}{\IfNoValueTF{#2}{\origcitet[#1]{#3}}{\origcitet[#1][#2]{#3}}}}}

% --- Sectioning and running heads ----------------------------------
\usepackage{titlesec}
\titleformat{\section}
  {\Large\bfseries}{\thesection}{0.6em}{}
  [\vspace{-0.5em}{\color{rulecol}\rule{\textwidth}{0.8pt}}]
\titlespacing*{\section}{0pt}{1.6em}{0.8em}
% \paragraph is the fourth level; make it a display heading rather than
% a run-in one. It is left unnumbered (see secnumdepth below).
% \titleformat{\paragraph}[block]
%   {\normalsize\bfseries}{\theparagraph}{0.6em}{}
% \titlespacing*{\paragraph}{0pt}{1.2em}{0.4em}

\usepackage{fancyhdr}
\pagestyle{fancy}
\fancyhf{}
\fancyhead[L]{\small\itshape\projectname\ --- \doctitleshort}
\fancyhead[R]{\small\thepage}
\fancyfoot[L]{\small\color{placeholdercol}\docversion}
\fancyfoot[R]{\small\color{placeholdercol}\docstatus}
\renewcommand{\headrulewidth}{0.4pt}
\renewcommand{\footrulewidth}{0.4pt}
% Plain pages (first page of the ToC, etc.) keep the same footer.
\fancypagestyle{plain}{%
  \fancyhf{}%
  \fancyfoot[L]{\small\color{placeholdercol}\docversion}%
  \fancyfoot[R]{\small\color{placeholdercol}\docstatus}%
  \renewcommand{\headrulewidth}{0pt}%
  \renewcommand{\footrulewidth}{0.4pt}%
}

% --- Spacing and table helpers -------------------------------------
\setlength{\parskip}{0.45em}
\setlength{\parindent}{0pt}
\setlength{\headheight}{14pt}
\newcolumntype{L}{>{\raggedright\arraybackslash}X}          % tabularx only

\newcolumntype{P}[1]{>{\raggedright\arraybackslash}p{#1}}   % longtable / tabular

% --- Numbering depth -----------------------------------------------
% Numbering stops at \subsubsection, so \paragraph headings carry no
% number. tocdepth 3 lists down to \subsubsection (a clustering method)
% but not the \paragraph headings beneath it.
\setcounter{secnumdepth}{3}
\setcounter{tocdepth}{3}

% --- Drafting helpers ----------------------------------------------
% \placeholder{...} : grey guidance text; delete as real content lands.
% \todo{...}        : visible red marker for outstanding work.
% Set \draftfalse (see main.tex) to hide both in a final build.
\newif\ifdraft
% \texorpdfstring keeps them usable inside section titles / bookmarks.
\newcommand{\placeholder}[1]{\texorpdfstring{\ifdraft{\color{placeholdercol}\itshape #1}\fi}{#1}}
\newcommand{\todo}[1]{\texorpdfstring{\ifdraft{\color{red}\bfseries [TODO: #1]}\fi}{[TODO: #1]}}

% --- Shorthand macros used across the technique sections -----------
\newcommand{\technique}[1]{\textsc{#1}}   % algorithm names in running text
